\section{投资问题：扩产是价值期权还是融资负担}

文莱二期规划1,200万吨／年，目标在2028年底建成；建成后文莱炼厂合计能力将达到2,000万吨／年。项目能够扩大重油加工、柴油、PX、苯和聚丙烯等产品组合，但当前披露主要是建设启动、税惠批文、贷款意向函和股东贷款承诺函。核心判断是，二期具有产业协同价值，却尚不具备进入基准目标价所需的融资、回报和建设曲线证据；主模型对二期维持零基准信用，只在牛情景和事件监控中体现。

公司公告称已取得文莱伊斯兰银行融资贷款意向函及达迈控股股东贷款承诺函，并表述项目建设所需资金已筹措到位；同时公告风险提示明确，仍存在建设资金未及时筹措到位、无法按期建成、市场和价格变化影响经济效益的风险。本报告据此区分“融资支持文件已取得”与“完整金融交割已完成”：在贷款金额、利率、期限、提款条件、资本金和担保结构未完整披露前，不能把意向函等同于已到账资金。

\begin{exhibitbox}[文莱二期项目状态]
\noindent\begin{tabularx}{\textwidth}{L{3.25cm}L{2.45cm}L{2.55cm}X}
\toprule
\textbf{项目字段} & \textbf{已披露} & \textbf{当前判断} & \textbf{估值门槛} \\
\midrule
设计产能 & 1,200万吨／年 & 规划明确 & 不等于实现销量 \\
产品结构 & 柴油、PX、苯、聚丙烯等 & 重油加工与高附加值方向 & 需产品收率和市场验证 \\
目标工期 & 力争2028年底建成 & 管理层目标，不是保证 & 按季度建设曲线核验 \\
2025年末进度 & 3.17\% & 早期建设 & 基准不给完工价值 \\
2026年3月末已投资 & 4.34亿美元，主要为土地平整 & 前期投入已发生 & 需后续资本金和工程包 \\
融资文件 & 税惠批文、贷款意向、股东贷款承诺 & 支持方向成立 & 金额、利率、期限、提款条件和到账 \\
一期协同 & 合计能力2,000万吨／年 & 规模与产品协同可能 & 需项目IRR、增量净利与现金回流 \\
\bottomrule
\end{tabularx}
\end{exhibitbox}
\sourcenote{HY-OFF-P2（2026-01-05）；HY-OFF-2025AR（进度截至2025-12-31）；HY-OFF-RATING（已投资截至2026-03-31）。贷款意向函和承诺函不视为完整金融交割。}

二期的估值含义不是简单“产能增加150\%”。项目建成前消耗资本，建成后仍需原油供给、市场、负荷、税、折旧、利息和少数股东共同决定回报。若以峰值裂解价差评估二期，资本开支与周期顶部会被同时低估。

\section{资本开支不是只有文莱二期}

截至2026年3月末，评级材料列示三项主要国内化纤在建项目计划投资合计173.97亿元、已投资134.92亿元，包括杭州逸通140万吨功能性纤维升级、宿迁逸达110万吨差别化纤维以及广西己内酰胺--聚酰胺一期。部分产能已经投产或试生产，但剩余投入、爬坡和营运资金仍会影响现金。

此外，公司2026年5月公告拟建设240万吨高品质纤维用煤制乙二醇项目，总投资预计257亿元、建设周期约三年，资金来自自有及自筹；另有循环新材料等项目。煤制乙二醇可提高聚酯原料自给，却也显著扩大资本需求、煤化工执行和环保风险。文莱二期、煤制乙二醇和国内技改不能被分别用“项目融资”解释后忽略集团总融资约束。

\begin{exhibitbox}[主要资本项目与现金约束]
\small
\noindent\begin{tabularx}{\textwidth}{L{3.05cm}R{2.0cm}R{2.0cm}L{2.15cm}X}
\toprule
\textbf{项目} & \textbf{计划投资} & \textbf{已投资} & \textbf{阶段} & \textbf{研究处理} \\
\midrule
杭州逸通140万吨升级 & 30.00亿元 & 13.16亿元 & 一期40万吨已投产，其余择机 & 剩余投入与市场匹配 \\
宿迁逸达110万吨差别化纤维 & 38.50亿元 & 18.21亿元 & 分期投产 & 关注库存与差异化毛利 \\
广西CPL--PA6一期 & 105.47亿元 & 103.55亿元 & 2025年10月试生产 & 关注爬坡现金与良率 \\
三项国内项目合计 & 173.97亿元 & 134.92亿元 & 建设／爬坡 & 已投入不等于稳定回报 \\
文莱二期 & 总投资未在本案完整披露 & 4.34亿美元 & 早期土地与建设 & 基准零信用，等金融交割 \\
煤制乙二醇240万吨 & 预计257亿元 & 未在本案形成统一进度 & 拟建、三年投入 & 单列重大资金与执行风险 \\
\bottomrule
\end{tabularx}
\end{exhibitbox}
\sourcenote{HY-OFF-RATING，数据截至2026-03-31；煤制乙二醇项目为评级材料引用公司2026年5月公告。文莱二期总投资未在本案当前公告中形成统一可审计数值。}

投资上，扩产选择权需要与资产负债表容量配对。2026Q1粗略净债务代理约574亿元、OCF为-25.73亿元；即使H1利润很高，也不能假设所有项目都能由内部现金无压力覆盖。若项目持续推进而经营现金没有转正，融资成本和稀释风险会提高。

\section{从项目进度到股东价值的五道门}

第一道门是工程：设计、采购、施工和试车是否按计划。第二道门是融资：贷款与资本金是否按明确成本和期限到账。第三道门是市场：产品结构是否匹配东南亚供需。第四道门是项目财务：负荷、毛利、税率和折旧后能否覆盖利息。第五道门是归母现金：70\%权益和担保结构下，现金能否回到上市公司。

\begin{exhibitbox}[二期估值升级阶梯]
\noindent\begin{tabularx}{\textwidth}{L{2.1cm}L{3.0cm}L{3.2cm}X}
\toprule
\textbf{阶段} & \textbf{必要证据} & \textbf{当前状态} & \textbf{估值信用} \\
\midrule
规划 & 实施协议、产能、产品 & 已具备 & 只确认战略方向 \\
融资 & 贷款金额、利率、期限、提款、资本金 & 意向／承诺，字段不全 & 基准为零 \\
建设 & 工程包、进度、累计投入 & 早期、已投资4.34亿美元 & 事件期权 \\
投产 & 试车、收率、负荷、认证／销售 & 尚未发生 & 不进入盈利预测 \\
现金回报 & 税后利润、少数股东、偿债与分红 & 尚未发生 & 满足后进入SOTP／FCFE \\
\bottomrule
\end{tabularx}
\end{exhibitbox}
\sourcenote{HY-OFF-P2；analysis/coverage\_gap\_matrix.md；analysis/value\_chain\_economics.md，研究截止2026-07-22。}

\begin{keyinsight}[资本开支的“所以呢”]
二期越有战略价值，越需要资本纪律证明。\textbf{行动上，目标价不提前计入二期；只有完整金融交割、可核验建设曲线和明确回报门槛出现后，才把项目从尾部期权纳入基准。}
\end{keyinsight}

\begin{riskbox}[资本开支失效条件]
若文莱二期或煤制乙二醇项目在经营现金流未修复时加速投入，贷款利率、资本金和担保继续不透明，或建设进度落后于2028目标，则应提高资本成本、下调FCFE和基准PE。\textbf{行动上，不用项目“全面启动”公告抵消现金与杠杆恶化。}
\end{riskbox}
